% =============================================================================
% UBS COE EXPERIENCE DESIGN - FRAGE-ABO OFFERTE
% =============================================================================
% FehrAdvice & Partners AG
% Compile with: xelatex UBS_037_frage_abo_coe_xd_proposal.tex
% =============================================================================

\documentclass{fehradvice-report}

% =============================================================================
% DOCUMENT METADATA
% =============================================================================

\title{Verhaltensökonomisches Frage-Abo}
\subtitle{für das Center of Excellence Experience Design}
\author{FehrAdvice \& Partners AG}
\date{Januar 2026}

\client{UBS Switzerland AG}
\project{CoE Experience Design -- Laufende Begleitung}
\version{1.0}
\classification{Vertraulich}

% =============================================================================
% DOCUMENT
% =============================================================================

\begin{document}

% Title page
\maketitle

% Table of contents
\tableofcontents
\newpage

% =============================================================================
\section{Executive Summary}
% =============================================================================

\subsection{Hintergrund und Problemstellung}

Das Center of Excellence for Experience Design (CoE XD) hat sich als zentrale Einheit für kundenzentrierte Lösungen bei UBS Switzerland etabliert. Mit dem erfolgreich abgeschlossenen Capability Building Projekt (UBS\_036) verfügt das Team nun über ein solides Fundament in verhaltensökonomischen Methoden, Storytelling und Experimentation.

Die Herausforderung besteht nun darin, dieses Wissen im Arbeitsalltag kontinuierlich anzuwenden und weiterzuentwickeln. Im täglichen Geschäft entstehen laufend Fragestellungen zu Customer Journeys, Touchpoint-Optimierungen und Design-Entscheidungen, die von einer verhaltensökonomischen Perspektive profitieren würden.

\subsection{Projektziel und Ansatz}

\begin{fainfobox}[Ziel des Angebots]
Etablierung eines strukturierten Prozesses für monatliche verhaltensökonomische Analysen, der:
\begin{itemize}
    \item Die im Capability Building erlernten Methoden in der Praxis verankert
    \item Dem CoE XD Team einen kontinuierlichen Sparringspartner bietet
    \item Design-Entscheidungen auf eine wissenschaftlich fundierte Basis stellt
    \item Die strategische Positionierung des CoE XD innerhalb von UBS stärkt
\end{itemize}
\end{fainfobox}

\subsection{Erwartete Ergebnisse und Mehrwerte}

\begin{itemize}
    \item \textbf{Für das CoE XD Team:} Schnelle, fundierte Antworten auf Experience-Design-Fragen -- von Journey-Optimierungen bis zu Touchpoint-Entscheidungen
    \item \textbf{Für UBS Kund:innen:} Bessere Kundenerlebnisse durch evidenzbasiertes Design
    \item \textbf{Für UBS Switzerland:} Stärkung des CoE XD als strategischer Partner mit messbarem Impact
\end{itemize}

\newpage
% =============================================================================
\section{Ausgangslage und Herausforderung}
% =============================================================================

\subsection{Aufbauend auf dem Capability Building}

Mit dem Projekt UBS\_036 hat das CoE XD wichtige Grundlagen geschaffen:

\begin{itemize}
    \item Ein gemeinsames Verständnis verhaltensökonomischer Prinzipien
    \item Erste BE Ambassadors, die als interne Multiplikatoren wirken
    \item Storytelling-Fähigkeiten zur Kommunikation von Design-Entscheidungen
    \item Ein Fundament für das langfristige Curriculum
\end{itemize}

\begin{fainfobox}[Die nächste Stufe]
Diese Fähigkeiten müssen nun im Alltag angewendet und vertieft werden. Dafür braucht es einen strukturierten Rahmen, der sicherstellt, dass verhaltensökonomische Analysen nicht ad hoc, sondern systematisch erfolgen.
\end{fainfobox}

\subsection{Warum ein strukturierter Analyseprozess notwendig ist}

Im Arbeitsalltag des CoE XD entstehen laufend Fragestellungen, die von einer verhaltensökonomischen Perspektive profitieren würden:

\begin{center}
\begin{tabular}{p{4cm}p{9cm}}
\toprule
\rowcolor{fadarkblue}
\fahead{Fragetyp} & \fahead{Beispiele} \\
\midrule
\rowcolor{fawhite}
Journey-Design & Wo in der Customer Journey verlieren wir Kund:innen? Welche Friction Points sind kritisch? \\
\rowcolor{falightgray}
Touchpoint-Optimierung & Wie sollte ein Onboarding-Flow gestaltet sein? Welche Defaults sind optimal? \\
\rowcolor{fawhite}
Choice Architecture & Wie viele Optionen sind zu viel? Wie präsentieren wir Produktauswahl? \\
\rowcolor{falightgray}
Behavioral Barriers & Warum nutzen Kund:innen Feature X nicht? Was sind die versteckten Barrieren? \\
\rowcolor{fawhite}
A/B-Test Design & Welche Varianten sollten wir testen? Wie interpretieren wir Ergebnisse verhaltensökonomisch? \\
\bottomrule
\end{tabular}
\end{center}

\vspace{0.5cm}

Eine systematische verhaltensökonomische Analyse kann:

\begin{itemize}
    \item \textbf{Versteckte Barrieren aufdecken}, die in traditionellen UX-Analysen übersehen werden
    \item \textbf{Wirkungsprognosen liefern}, bevor Ressourcen in die Umsetzung fliessen
    \item \textbf{Design-Entscheidungen absichern} mit wissenschaftlicher Fundierung
    \item \textbf{Interne Akzeptanz stärken}, weil Empfehlungen nachvollziehbar begründet sind
\end{itemize}

\subsection{Bedeutung für das CoE XD}

Mit der Etablierung eines strukturierten Prozesses für monatliche verhaltensökonomische Analysen:

\begin{itemize}
    \item Werden die im Capability Building erlernten Methoden \textbf{kontinuierlich angewendet und vertieft}
    \item Erhält das CoE XD eine \textbf{robuste Entscheidungsgrundlage} für Design-Empfehlungen
    \item Stärkt das Team seine \textbf{Positionierung als strategischer Partner} mit evidenzbasiertem Ansatz
    \item Entsteht ein \textbf{Lernprozess}, der das kollektive Wissen über Zeit aufbaut
\end{itemize}

FehrAdvice \& Partners agiert dabei als Sparringspartner für das CoE XD -- analysierend, reflektierend, verhaltensökonomisch fundiert und auf die spezifischen Herausforderungen von Experience Design zugeschnitten.

\newpage
% =============================================================================
\section{Zielsetzung}
% =============================================================================

Das Ziel dieses Projekts ist es, einen strukturierten Prozess für monatliche verhaltensökonomische Analysen zu etablieren, der die Design-Entscheidungen des CoE XD auf eine wissenschaftliche Basis stellt und die strategische Positionierung des Teams innerhalb von UBS stärkt.

\subsection{Strategische Ziele}

\begin{itemize}
    \item[$\checkmark$] \textbf{Capability Building fortführen:} Die im Projekt UBS\_036 erlernten Methoden werden im Alltag verankert und kontinuierlich weiterentwickelt.
    \item[$\checkmark$] \textbf{Strategische Positionierung stärken:} Das CoE XD etabliert sich als Quelle evidenzbasierter Design-Empfehlungen -- nicht nur als Ausführungseinheit.
    \item[$\checkmark$] \textbf{Messbare Wirkung nachweisen:} Durch dokumentierte Analysen und deren Umsetzungserfolge wird der Wertbeitrag des CoE XD sichtbar.
\end{itemize}

\subsection{Verhaltensorientierte Ziele}

\begin{itemize}
    \item[$\checkmark$] \textbf{Wissenschaftlich fundierte Design-Entscheidungen:} Fragestellungen werden nicht mehr nur auf Basis von Best Practices oder Intuition beantwortet, sondern mit verhaltensökonomischen Modellen analysiert.
    \item[$\checkmark$] \textbf{Behavioral Barriers identifizieren:} Analysen zeigen, welche versteckten Barrieren Kund:innen an der Nutzung von UBS-Services hindern -- und wie diese überwunden werden können.
    \item[$\checkmark$] \textbf{Konsistenz durch systematische Methodik:} Die Analysen beruhen auf dem Behavioral Change Model und dem EBF Framework, was Transparenz und Nachvollziehbarkeit sichert.
\end{itemize}

\subsection{Operative Ziele}

\begin{itemize}
    \item[$\checkmark$] \textbf{Pragmatische Analysen für den Alltag:} Jede Fragestellung führt zu konkreten Design-Empfehlungen -- direkt umsetzbar für Journeys, Touchpoints oder Features.
    \item[$\checkmark$] \textbf{Frühwarnsystem für UX-Risiken:} Potenzielle Friction Points oder Adoption-Barrieren werden identifiziert, bevor sie in der Umsetzung zu Problemen führen.
    \item[$\checkmark$] \textbf{Kontinuierliches Lernen etablieren:} Durch den monatlichen Rhythmus entsteht eine Wissensbasis, die über Zeit wächst und dem gesamten Team zur Verfügung steht.
\end{itemize}

\newpage
% =============================================================================
\section{Unser Vorschlag zum Vorgehen}
% =============================================================================

Das Vorgehen folgt einem klar strukturierten Prozess, der sicherstellt, dass jede Fragestellung konsistent, wissenschaftlich fundiert und praxisnah bearbeitet wird.

\subsection{Phase 1: Definition der Fragestellungen}

\begin{itemize}
    \item Das CoE XD definiert monatlich die relevanten inhaltlichen Fragestellungen, die im Arbeitsalltag entstehen. Diese können Customer Journeys, Touchpoints, Features oder strategische Design-Entscheidungen betreffen.
    \item Die Fragestellungen werden an FehrAdvice übermittelt, zusammen mit allen notwendigen Unterlagen (z.B. Journey Maps, Prototypen, Analytics-Daten, User Research).
    \item Zusätzlich wird festgelegt, ob die Bearbeitung als \textbf{Kurzanalyse} (3--5 Seiten) oder als \textbf{umfassende Analyse} (8--10 Seiten) erfolgen soll.
\end{itemize}

\subsection{Phase 2: Verhaltensökonomische Analyse}

Auf Basis der eingereichten Materialien erfolgt eine systematische Analyse entlang verhaltensökonomischer Logik:

\begin{itemize}
    \item \textbf{Kontextfaktoren ($\Psi$-Dimensionen):} Welche situativen Faktoren beeinflussen das Kundenverhalten?
    \item \textbf{Behavioral Barriers:} Welche kognitiven, emotionalen oder praktischen Hürden existieren?
    \item \textbf{Nutzenfunktionen:} Welche Utility-Dimensionen (FEPSDE) sind für die Kund:innen relevant?
    \item \textbf{Intervention Options:} Welche BE-Tools (Defaults, Framing, Salience, etc.) sind geeignet?
\end{itemize}

Die Analyse ist wissenschaftlich fundiert und basiert auf dem EBF Framework sowie dem Behavioral Change Model. Aus den Analysen werden \textbf{praxisnahe Design-Empfehlungen} abgeleitet.

\subsection{Phase 3: Aufbereitung der Ergebnisse}

Die Ergebnisse werden in Form eines klar strukturierten Assessments dokumentiert:

\begin{itemize}
    \item \textbf{Kurzanalyse} (3--5 Seiten) für fokussierte, kleinere Fragestellungen
    \item \textbf{Umfassende Analyse} (8--10 Seiten) für komplexere, strategisch bedeutende Fragestellungen
\end{itemize}

\begin{fainfobox}[Einheitliche Berichtsstruktur]
\begin{enumerate}
    \item \textbf{Ausgangslage} -- Kontext der Design-Herausforderung
    \item \textbf{Fragestellung} -- Präzise Formulierung des zu analysierenden Problems
    \item \textbf{Verhaltensökonomische Analyse} -- Detaillierte Untersuchung mit BE-Framework
    \item \textbf{Design-Empfehlungen} -- Konkrete, umsetzbare Handlungsoptionen
\end{enumerate}
\end{fainfobox}

\newpage
% =============================================================================
\section{Unser Leistungsangebot}
% =============================================================================

\subsection{Übersicht und Abgrenzung der Leistungsmodule}

Die drei Module unterscheiden sich primär im Umfang und der Frequenz der verhaltensökonomischen Analysen.

\begin{center}
\begin{tabular}{p{2.5cm}p{5.5cm}p{5cm}}
\toprule
\rowcolor{fadarkblue}
\fahead{Modul} & \fahead{Eignung} & \fahead{Umfang pro Monat} \\
\midrule
\rowcolor{fawhite}
\textbf{Small} & Gezielte Bearbeitung einzelner Design-Fragestellungen & 1 umfassend ODER 2 kurz \\
\rowcolor{falightgray}
\textbf{Medium} & Mehrere Fragestellungen, Zusammenhänge sichtbar & 2 umfassend ODER 3 kurz \\
\rowcolor{fawhite}
\textbf{Large} & Maximale Flexibilität, laufende Begleitung & Nach Absprache \\
\bottomrule
\end{tabular}
\end{center}

\vspace{1cm}

% -----------------------------------------------------------------------------
\subsection{Modul Small -- Einzelne Design-Fragestellungen präzise beantworten}
% -----------------------------------------------------------------------------

\textbf{Ziele:}
\begin{itemize}
    \item[$\checkmark$] \textbf{Klarheit schaffen:} Antworten auf präzise Fragestellungen aus dem CoE XD Arbeitsalltag -- schnell, nachvollziehbar, nutzbar.
    \item[$\checkmark$] \textbf{Design-Entscheidungen absichern:} Optionen aufzeigen, damit Entscheidungen nicht auf Intuition allein beruhen.
    \item[$\checkmark$] \textbf{Risiken erkennen:} Frühzeitige Hinweise auf potenzielle UX-Probleme oder Adoption-Barrieren.
\end{itemize}

\textbf{Leistungen:}
\begin{itemize}
    \item Systematische Analyse von \textbf{1 umfassenden Analyse} (8--10 Seiten) ODER \textbf{2 Kurzanalysen} (je 3--5 Seiten) pro Monat
    \item 3--4 Wochen pro Analyse, abhängig von Umfang und Materiallage
    \item Auswertung der von UBS bereitgestellten Unterlagen (Journey Maps, Prototypen, Analytics, User Research)
    \item Wissenschaftlich fundierte Bearbeitung mit Ableitung praxisnaher Design-Empfehlungen
\end{itemize}

\textbf{Out of scope:}
\begin{itemize}
    \item[$\times$] Empirische Primärerhebung (quantitativ, qualitativ, experimentell)
    \item[$\times$] Operative Umsetzung oder Design-Implementierung
    \item[$\times$] Entwicklung von Prototypen oder UI-Designs
\end{itemize}

\textbf{Lieferobjekte:}\\
Bericht in Form von 2 Kurzanalysen (3--5 Seiten) oder 1 umfassenden Analyse (8--10 Seiten) pro Monat mit einheitlicher Struktur: Ausgangslage -- Fragestellung -- Verhaltensökonomische Analyse -- Design-Empfehlungen

\vspace{0.3cm}
\begin{center}
\fcolorbox{fadarkblue}{falightgray}{%
\parbox{0.9\textwidth}{%
\centering
\textbf{Dauer:} 3 Monate (d.h. drei Monatszyklen von Analysen)\\[0.3cm]
\textbf{\Large Preis: CHF 12'000.-}
}%
}
\end{center}

\newpage
% -----------------------------------------------------------------------------
\subsection{Modul Medium -- Vertiefte Analysen mehrerer Design-Fragestellungen}
% -----------------------------------------------------------------------------

\textbf{Ziele:}
\begin{itemize}
    \item[$\checkmark$] \textbf{Themen verbinden:} Analysen mehrerer Fragen so aufbereiten, dass Zusammenhänge und Patterns sichtbar werden.
    \item[$\checkmark$] \textbf{Prioritäten setzen:} Das CoE XD erhält eine fundierte Basis, um Design-Ressourcen gezielt zu steuern.
    \item[$\checkmark$] \textbf{Design-Optionen vergleichen:} Alternativen nachvollziehbar nebeneinanderstellen für bessere Entscheidungen.
\end{itemize}

\textbf{Leistungen:}
\begin{itemize}
    \item Systematische Analyse von \textbf{2 umfassenden Analysen} (8--10 Seiten) ODER \textbf{3 Kurzanalysen} (je 3--5 Seiten) pro Monat
    \item 3--4 Wochen pro Analyse, abhängig von Umfang und Materiallage
    \item Auswertung der von UBS bereitgestellten Unterlagen
    \item Wissenschaftlich fundierte Bearbeitung mit Ableitung praxisnaher Design-Empfehlungen
\end{itemize}

\textbf{Out of scope:}
\begin{itemize}
    \item[$\times$] Empirische Primärerhebung (quantitativ, qualitativ, experimentell)
    \item[$\times$] Operative Umsetzung oder Design-Implementierung
    \item[$\times$] Entwicklung von Prototypen oder UI-Designs
\end{itemize}

\textbf{Lieferobjekte:}\\
Bericht in Form von 3 Kurzanalysen (3--5 Seiten) oder 2 umfassenden Analysen (8--10 Seiten) pro Monat mit einheitlicher Struktur.

\vspace{0.3cm}
\begin{center}
\fcolorbox{fadarkblue}{falightgray}{%
\parbox{0.9\textwidth}{%
\centering
\textbf{Dauer:} 3 Monate (d.h. drei Monatszyklen von Analysen)\\[0.3cm]
\textbf{\Large Preis: CHF 18'000.-}
}%
}
\end{center}

\vspace{1cm}
% -----------------------------------------------------------------------------
\subsection{Modul Large -- Laufende Analyse und Begleitung}
% -----------------------------------------------------------------------------

\textbf{Ziele:}
\begin{itemize}
    \item[$\checkmark$] \textbf{Flexibel reagieren:} Auch kurzfristige, unerwartete Fragestellungen können fundiert beantwortet werden.
    \item[$\checkmark$] \textbf{Laufend begleitet:} Das CoE XD verfügt über eine kontinuierliche Reflexionsbasis für Design-Entscheidungen.
    \item[$\checkmark$] \textbf{Strategische Wirkung absichern:} Design-Empfehlungen werden regelmässig validiert, um Kundenzufriedenheit und Business Impact zu sichern.
\end{itemize}

\textbf{Leistungen:}
\begin{itemize}
    \item Laufende Bearbeitung von Fragestellungen in flexibler Tiefe und Frequenz
    \item Möglichkeit zur kurzfristigen Bearbeitung von ad hoc Themen
    \item Kontinuierliche Begleitung als Sparringspartner
    \item Teilnahme an ausgewählten CoE XD Meetings (nach Absprache)
\end{itemize}

\textbf{Lieferobjekte:}
\begin{itemize}
    \item Laufende Assessments in vereinbarter Tiefe
    \item Flexible Reports (3--10 Seiten je nach Fragestellung)
    \item Periodische Executive Summaries mit übergreifenden Erkenntnissen
    \item Quartalsweise Review der wichtigsten Learnings
\end{itemize}

\vspace{0.3cm}
\begin{center}
\fcolorbox{fadarkblue}{falightgray}{%
\parbox{0.9\textwidth}{%
\centering
\textbf{Dauer:} Fortlaufend, Umfang und Rhythmus nach Absprache\\[0.3cm]
\textbf{\Large Preis: nach Absprache}
}%
}
\end{center}

\newpage
% =============================================================================
\subsection{Leistungsübersicht}
% =============================================================================

\begin{center}
\begin{tabular}{p{7cm}ccc}
\toprule
\rowcolor{fadarkblue}
\fahead{Leistungsbausteine} & \fahead{Modul S} & \fahead{Modul M} & \fahead{Modul L} \\
\midrule
\rowcolor{fawhite}
Kick-off \& Zielklärung & $\checkmark$ & $\checkmark$ & $\checkmark$ \\
\rowcolor{falightgray}
Auswertung UBS-Unterlagen (Journeys, Prototypen, Analytics) & $\checkmark$ & $\checkmark$ & $\checkmark$ \\
\rowcolor{fawhite}
Wissenschaftlich fundierte Analyse \& Design-Empfehlungen & $\checkmark$ & $\checkmark$ & $\checkmark$ \\
\rowcolor{falightgray}
Strukturierter Analysebericht (kurz oder umfassend) & $\checkmark$ & $\checkmark$ & $\checkmark$ \\
\rowcolor{fawhite}
Standardumfang (2 kurz oder 1 umfassend) & $\checkmark$ & & \\
\rowcolor{falightgray}
Erweiterter Umfang (3 kurz oder 2 umfassend) & & $\checkmark$ & \\
\rowcolor{fawhite}
Flexibilität für ad hoc Fragestellungen & & & $\checkmark$ \\
\rowcolor{falightgray}
Periodische Executive Summaries & & & $\checkmark$ \\
\rowcolor{fawhite}
Laufende Begleitung als Sparringspartner & & & $\checkmark$ \\
\midrule
\rowcolor{falightgray}
\textbf{Dauer} & \textbf{3 Monate} & \textbf{3 Monate} & \textbf{Fortlaufend} \\
\rowcolor{fadarkblue}
\fahead{Preis} & \fahead{CHF 12'000.-} & \fahead{CHF 18'000.-} & \fahead{n. Abspr.} \\
\bottomrule
\end{tabular}
\end{center}

\newpage
% =============================================================================
\section{Der Nutzen für das CoE Experience Design}
% =============================================================================

Das Angebot liefert für das CoE XD einen unmittelbaren und langfristigen Mehrwert.

\subsection{Konkrete Unterstützung im Arbeitsalltag}

\begin{itemize}
    \item[$\checkmark$] \textbf{Schnelle Orientierung:} Klare, wissenschaftlich fundierte Antworten auf Design-Fragen -- sei es zu Journeys, Touchpoints oder Features.
    \item[$\checkmark$] \textbf{Reduktion von Unsicherheit:} Design-Entscheidungen stützen sich auf transparente Analysen statt nur auf Best Practices.
    \item[$\checkmark$] \textbf{Zeitersparnis:} Statt lange interne Diskussionen zu führen, kann auf präzise Berichte mit Design-Empfehlungen zurückgegriffen werden.
\end{itemize}

\subsection{Strategischer Mehrwert für das CoE XD}

\begin{itemize}
    \item[$\checkmark$] \textbf{Positionierung stärken:} Das CoE XD etabliert sich als Quelle evidenzbasierter Empfehlungen -- ein Differenzierungsmerkmal innerhalb UBS.
    \item[$\checkmark$] \textbf{Capability Building fortführen:} Die im Projekt UBS\_036 erlernten Methoden werden kontinuierlich angewendet und vertieft.
    \item[$\checkmark$] \textbf{Messbare Wirkung nachweisen:} Dokumentierte Analysen und deren Umsetzungserfolge machen den Wertbeitrag des CoE XD sichtbar.
\end{itemize}

\subsection{Nutzen für UBS Kund:innen}

\begin{itemize}
    \item[$\checkmark$] \textbf{Bessere Kundenerlebnisse:} Design-Entscheidungen basieren auf einem tiefen Verständnis von Kundenverhalten.
    \item[$\checkmark$] \textbf{Weniger Friction:} Behavioral Barriers werden identifiziert und systematisch abgebaut.
    \item[$\checkmark$] \textbf{Höhere Adoption:} Neue Features und Services werden so gestaltet, dass sie tatsächlich genutzt werden.
\end{itemize}

\begin{fasuccessbox}[Zusammenfassung]
UBS Switzerland gewinnt mit diesem Angebot ein Werkzeug, das den Arbeitsalltag des CoE Experience Design erleichtert, die strategische Positionierung des Teams stärkt und evidenzbasiertes Design als Standard etabliert.
\end{fasuccessbox}

\newpage
% =============================================================================
\appendix
\section{Verbindung zum Capability Building (UBS\_036)}
% =============================================================================

Dieses Angebot baut direkt auf dem erfolgreich abgeschlossenen Capability Building Projekt auf:

\begin{center}
\begin{tabular}{p{6cm}p{6cm}}
\toprule
\rowcolor{fadarkblue}
\fahead{UBS\_036 (abgeschlossen)} & \fahead{Frage-Abo (dieses Angebot)} \\
\midrule
\rowcolor{fawhite}
Grundlagen vermittelt & Anwendung im Alltag \\
\rowcolor{falightgray}
BE-Sprache etabliert & BE-Analysen durchgeführt \\
\rowcolor{fawhite}
Ambassadors ausgebildet & Ambassadors unterstützt \\
\rowcolor{falightgray}
Experimentation Mindset & Experimente begleitet \\
\bottomrule
\end{tabular}
\end{center}

% =============================================================================
\section{Beispiele für typische Fragestellungen}
% =============================================================================

\textbf{Journey-Design:}
\begin{itemize}
    \item «Wo in der Onboarding-Journey verlieren wir die meisten Neukund:innen?»
    \item «Wie sollte der Wechsel von App zu Filiale gestaltet sein?»
\end{itemize}

\textbf{Touchpoint-Optimierung:}
\begin{itemize}
    \item «Welche Defaults sollten wir im Anlageberater setzen?»
    \item «Wie reduzieren wir die Komplexität bei der Produktauswahl?»
\end{itemize}

\textbf{Feature Adoption:}
\begin{itemize}
    \item «Warum wird Feature X kaum genutzt, obwohl es gut bewertet wird?»
    \item «Wie können wir die Nutzung der neuen Budgeting-Funktion steigern?»
\end{itemize}

\textbf{A/B-Test Design:}
\begin{itemize}
    \item «Welche Varianten sollten wir für den neuen Checkout-Flow testen?»
    \item «Wie interpretieren wir die Ergebnisse aus verhaltensökonomischer Sicht?»
\end{itemize}

% =============================================================================
\section{Über FehrAdvice \& Partners}
% =============================================================================

FehrAdvice \& Partners ist das führende Beratungsunternehmen für Behavioral Economics im deutschsprachigen Raum. Gegründet von Prof. Ernst Fehr (Universität Zürich), verbinden wir wissenschaftliche Exzellenz mit praktischer Anwendung.

\textbf{Unsere Expertise:}
\begin{itemize}
    \item 15+ Jahre Erfahrung in Behavioral Economics Beratung
    \item 200+ Projekte in verschiedenen Branchen
    \item Wissenschaftliches Fundament durch Prof. Ernst Fehr
\end{itemize}

\vfill

\begin{center}
\textcolor{fadarkblue}{\rule{0.8\textwidth}{1pt}}\\[0.5cm]
\textbf{FehrAdvice \& Partners AG}\\
Klausstrasse 4 | 8008 Zürich\\
www.fehradvice.com
\end{center}

\end{document}
